\documentclass[../../../include/open-logic-section]{subfiles}

\begin{document}

\olfileid{sth}{choice}{woproblem}
\olsection{مسئلهٔ خوش‌ترتیبی}

آشکار است که پذیرفتن یا نپذیرفتن خوش‌ترتیبی پیامدهای فراوانی دارد.
اما بحث دربارهٔ این اصل بیشتر بر اصلی هم‌ارز، یعنی اصل موضوع انتخاب،
متمرکز بوده است. پس اکنون به آن می‌پردازیم (و هم‌ارزی را نیز اثبات
می‌کنیم).

کانتور در \citeyear{Cantor1883} پشتیبانی خود را از اصل موضوع
خوش‌ترتیبی ابراز کرد و آن را «قانونی برای اندیشه که به نظرم بنیادی،
سرشار از پیامد و به‌ویژه به سبب اعتبار عامش درخور توجه است» نامید
(به نقل از \citeauthor{Potter2004}
\citeyear[p.~243]{Potter2004}). اما سرانجام کانتور قانع شد که این
«اصل موضوع» نیازمند اثبات است. جامعهٔ ریاضی نیز چنین اندیشید.

این مسئله را زرملو در \citeyear{Zermelo1904} «حل کرد». برای توضیح
راه‌حل او، به چند تعریف نیاز داریم.

\begin{defn}
تابع $f$ یک \emph{تابع انتخاب} است اگر و تنها اگر برای هر
$x \in \dom{f}$ داشته باشیم $f(x) \in x$. می‌گوییم $f$ یک
\emph{تابع انتخاب برای $A$} است اگر و تنها اگر $f$ تابع انتخابی با
$\dom{f} = A \setminus \{\emptyset\}$ باشد.
\end{defn}

از نظر شهودی، تابع انتخاب برای $A$ از هر مجموعهٔ (ناتهیِ)
$x \in A$ یک عضو مشخص، یعنی $f(x)$، را \emph{برمی‌گزیند}. اکنون اصل
موضوع انتخاب چنین است:

\begin{axiom}[انتخاب]
	هر مجموعه دارای یک تابع انتخاب است.
\end{axiom}

زرملو نشان داد که انتخاب مستلزم خوش‌ترتیبی است و برعکس:

\begin{thm}[در $\ZF$]\ollabel{thmwochoice}
خوش‌ترتیبی و انتخاب هم‌ارزند.
\end{thm}

\begin{proof}
\emph{از چپ به راست.} فرض کنید $A$ مجموعه‌ای از مجموعه‌ها باشد.
بنا بر اصل موضوع اجتماع، $\bigcup A$ وجود دارد؛ پس بنا بر
خوش‌ترتیبی، رابطه‌ای مانند $<$ وجود دارد که $\bigcup A$ را
خوش‌ترتیب می‌کند. اکنون بگذارید
$f(x) = \text{عضوِ $<$-کمینهٔ }x$. این تابعی انتخاب برای $A$ است.

\emph{از راست به چپ.} $A$ را ثابت کنید. اگر $A=\emptyset$، رابطهٔ
تهی آن را خوش‌ترتیب می‌کند؛ پس فرض کنید $A\neq\emptyset$. بنا بر
انتخاب، تابع انتخابی مانند $f$ برای
$\Pow{A} \setminus \{\emptyset\}$ وجود دارد. با استفاده از بازگشت
ترامتناهی، تابعی را تعریف کنید:
\begin{align*}
	g(0) &= f(A)\\
	g(\alpha) &=
		\begin{cases}
			\text{توقف!{}} &\text{اگر }A = \funimage{g}{\alpha}\\
			f(A \setminus \funimage{g}{\alpha}) & \text{در غیر این صورت}\\
		\end{cases}
\end{align*}
دستور «توقف!» تنها صورتی کوتاه برای تعریفی است که در غیر این صورت
طولانی‌تر می‌شد. یعنی هنگامی که برای نخستین بار
$A = \funimage{g}{\alpha}$، برای همهٔ
$\delta \geq \alpha$ بگذارید $g(\delta) = A$. در وهلهٔ نخست فقط
می‌توانیم مطمئن باشیم که این کار یک \emph{ترم} را تعریف می‌کند
(نگاه کنید به توضیحات پس از
\olref[sth][spine][recursion]{transrecursionschema})؛ اما نشان خواهیم
داد که بخش لازم از آن واقعاً یک تابع است.

چون $f$ تابع انتخاب است، برای هر $\alpha$ پیش از مرحلهٔ توقف داریم
$g(\alpha) =
f(A \setminus \funimage{g}{\alpha}) \in A \setminus
\funimage{g}{\alpha}$؛ یعنی
$g(\alpha) \notin \funimage{g}{\alpha}$. پس مقادیر $g$ پیش از توقف
!!{injective} هستند.

اکنون توجه کنید که ناگزیر متوقف می‌شویم؛ یعنی عدد ترتیبیِ (کمینهٔ)
$\alpha$ای وجود دارد که $A = g[\alpha]$. زیرا اگر هرگز متوقف
نمی‌شدیم، $g$ !!{injective} می‌بود و برای هر عدد ترتیبیِ $\alpha$
داشتیم $\cardless{\alpha}{A}$، برخلاف
\olref[hartogs]{HartogsLemma}. در کمینهٔ مرحلهٔ توقف، تحدیدِ $g$ به
دامنهٔ $\alpha$ بردی برابر با $A$ دارد.

با کنار هم گذاشتن این واقعیت‌ها، این تحدید یک !!a{bijection} از یک
عدد ترتیبی به $A$ است.
اکنون می‌توان از آن برای خوش‌ترتیب‌کردن $A$ استفاده کرد.
\end{proof}

پس خوش‌ترتیبی و انتخاب با هم پابرجا می‌مانند یا با هم فرومی‌ریزند.
اما پرسش همچنان باقی است: پابرجا می‌مانند یا فرومی‌ریزند؟

\end{document}
